\documentclass[12pt]{article}

\input{../../_assets/preambulo.tex}


\begin{document}

    % 1. Foto de fondo
    % 2. Título
    % 3. Encabezado Izquierdo
    % 4. Color de fondo
    % 5. Coord x del titulo
    % 6. Coord y del titulo
    % 7. Fecha

    
    \input{../../_assets/portada}
    \portadaExamen{ffccA4.jpg}{Ecuaciones\\Diferenciales II\\Examen XXII}{Ecuaciones Diferenciales II. Examen XXII}{MidnightBlue}{-8}{28}{2026}{}

    \begin{description}
        \item[Asignatura] Ecuaciones Diferenciales II.
        \item[Curso Académico] 2023/2024.
        \item[Grado] Doble Grado en Ingeniería Informática y Matemáticas.
        \item[Grupo] Único.
        \item[Profesor] Rafael Ortega Ríos.
        \item[Descripción] Convocatoria Ordinaria.
        \item[Fecha] 13 de junio de 2024.
        % \item[Duración] ---.
    
    \end{description}
    \newpage


    % ------------------------------------
    
    \begin{ejercicio}
        Encuentra el intervalo maximal de la solución del problema de valores iniciales
        \[
        \dot{x} = \frac{1}{t}x^2, \quad x(1) = 1.
        \]
    \end{ejercicio}

    \begin{ejercicio}
        Para cada $n \geq 1$ se considera el problema de valores iniciales
        \[
        \dot{x} = \cos x, \quad x(0) = \frac{1}{n}.
        \]
        \begin{enumerate}[label=\alph*)]
            \item Demuestra que existe una única solución $x_n(t)$ definida en toda la recta real.
            \item Demuestra que la sucesión de funciones $\{x_n\}_{n \geq 1}$ es equicontinua en el intervalo $\left]-\infty, +\infty\right[$.
        \end{enumerate}
    \end{ejercicio}

    \begin{ejercicio}
        Para cada $n \geq 1$ se considera el problema de valores iniciales
        \[
        \dot{x} = \sen x, \quad x(0) = \frac{1}{n}.
        \]
        Calcula, si existe, $\lim\limits_{n \to \infty} \int\limits_0^1 x_n(s)^2 ds$.
    \end{ejercicio}

    \begin{ejercicio}
        Encuentra todas las funciones continuas $x : \bb{R} \to \bb{R}$ que cumplen
        \[
        e^{x(t)} = \int_0^t e^{-x(s)} ds, \quad \forall t \in \bb{R}.
        \]
    \end{ejercicio}

    \begin{ejercicio}
        Construye una sucesión de funciones $\{f_n\}_{n \geq 0}$, $f_n : \left[-1, 1\right] \to \bb{R}$, que sea uniformemente acotada pero no sea equicontinua.
    \end{ejercicio}

    \begin{ejercicio}~
        \begin{enumerate}[label=\alph*)]
            \item Para cada $\veps \in \bb{R}$ la solución del problema de valores iniciales
            \[
            \ddot{x} + x + \veps^3 x^2 = 0, \quad x(0) = 0, \quad \dot{x}(0) = 1 + \veps
            \]
            se denota por $x(t; \veps)$. Calcula, si existe, $\deld{x}{\veps}(t; 0)$.
            \item Para cada $\veps \in \bb{R}$ la solución del problema de valores iniciales
            \[
            \ddot{x} + x + \veps^3 x^2 = 0, \quad x(\veps) = 0, \quad \dot{x}(\veps) = 1.
            \]
            se denota por $x(t; \veps)$. Calcula, si existe, $\deld{x}{\veps}(t; 0)$.
        \end{enumerate}
    \end{ejercicio}

    \newpage
    \setcounter{ejercicio}{0} % Reiniciar contador de ejercicios
    % \noindent
    % \textbf{Solución.}

\end{document}